% CALIDAD - diseño de 26 casos (CP-082..CP-107)
% GUÍA EQUIPO: Responsable asignado: José Mateo Auz (Calidad).
\paragraph{Módulo Calidad --- Ruteo, incidencias, concurrencia y rendimiento (26 casos: CP-082 a CP-107).}

\begin{disenio}
\caso{CP-082}{\textbf{Agrupación horaria 10:\allowbreak{}00+10:\allowbreak{}20 y separación 15:\allowbreak{}00.} NR-17, NR-18, NR-19; NP-07, NP-08 (C12.1, Alta)}{Pedidos de hoy: P1 (10:\allowbreak{}00), P2 (10:\allowbreak{}20) distantes y P3 (15:\allowbreak{}00) cercano a P1; margen configurado 30 min.}{1) Enviar GET /\allowbreak{}api/\allowbreak{}v1/\allowbreak{}route-proposals con rol ADMIN. 2) Revisar grupos propuestos y orden de paradas.}{Propuesta agrupa P1+P2 (10:\allowbreak{}00+10:\allowbreak{}20) en una ruta y separa P3 (15:\allowbreak{}00) en otra ruta sin mezclarlos por cercanía.}{P1 y P2 comparten ruta por ventana horaria compatible; P3 queda separado pese a menor distancia a P1.}
\caso{CP-083}{\textbf{Ruta válida con 1 parada sin exigir mínimo de 4.} NR-17, NR-18, NR-19; NP-07, NP-08 (C12.2, Alta)}{1 único pedido elegible de hoy con coordenadas viales válidas dentro de cobertura.}{1) Consultar GET /\allowbreak{}api/\allowbreak{}v1/\allowbreak{}route-proposals como ADMIN. 2) Verificar cantidad de rutas y paradas.}{Se genera 1 propuesta de ruta válida con 1 parada, origen en tienda y retorno, sin exigir mínimo de 4 pedidos.}{HTTP 200 con ruta propuesta de 1 parada utilizable y sin error de capacidad mínima.}
\caso{CP-084}{\textbf{Lote de 4 pedidos compatibles forma ruta de 4 paradas.} NR-17, NR-18, NR-19; NP-07, NP-08 (C12.3, Alta)}{4 pedidos elegibles de hoy en misma ventana horaria (ej. 10:\allowbreak{}00-10:\allowbreak{}25) y margen 30 min.}{1) Solicitar GET /\allowbreak{}api/\allowbreak{}v1/\allowbreak{}route-proposals como ADMIN. 2) Verificar número de rutas y paradas.}{El lote se agrupa en 1 ruta máxima de exactamente 4 paradas ordenadas cronológicamente.}{Total de paradas por ruta \ensuremath{\leq} 4; los 4 pedidos compatibles quedan incluidos en una sola ruta.}
\caso{CP-085}{\textbf{Lote de 5 pedidos dividido en máximo 4 paradas por ruta.} NR-17, NR-18, NR-19; NP-07, NP-08 (C12.4, Alta)}{5 pedidos elegibles de hoy en ventana compatible (margen 30 min) con coordenadas viales válidas.}{1) Invocar GET /\allowbreak{}api/\allowbreak{}v1/\allowbreak{}route-proposals como ADMIN. 2) Inspeccionar partición de rutas y paradas.}{El sistema divide los 5 pedidos en 2 rutas (4+1 o 3+2 según métricas viales), nunca una ruta de 5.}{Todas las rutas tienen entre 1 y 4 paradas; suma total de paradas igual a 5 sin duplicados.}
\caso{CP-086}{\textbf{Declarar emergencia pasa ruta a INTERRUPTED y bloquea.} NR-25; NP-11 (C19.1, Alta)}{Ruta IN\_PROGRESS con 1 entrega DELIVERED y 2 paradas pendientes; EMERGENCY\_DECLARE con routeId y requestId.}{1) Ejecutar POST /\allowbreak{}api/\allowbreak{}v1/\allowbreak{}dispatch/\allowbreak{}interventions con EMERGENCY\_DECLARE. 2) Consultar ruta, RouteUnitRecovery y pedidos.}{Ruta pasa a INTERRUPTED, conserva entrega realizada, crea RouteUnitRecovery por unidad pendiente y las bloquea.}{Estado INTERRUPTED, entrega previa intacta, un registro RouteUnitRecovery por unidad pendiente y bloqueo activo.}
\caso{CP-087}{\textbf{Retorno individual desbloquea solo unidad confirmada.} NR-25; NP-11 (C19.2, Alta)}{Dos unidades pendientes (U1, U2) en RouteUnitRecovery; Product.\allowbreak{}stock inicial S; UNIT\_RECOVERY\_CONFIRM para U1.}{1) Enviar UNIT\_RECOVERY\_CONFIRM con recoveryId de U1. 2) Consultar U1, U2 y Product.\allowbreak{}stock.}{U1 registra returnedAt y se desbloquea; U2 permanece bloqueada y Product.\allowbreak{}stock sigue igual a S.}{Solo U1 queda liberada; U2 conserva returnedAt NULL y no se genera movimiento de suma en Product.\allowbreak{}stock.}
\caso{CP-088}{\textbf{Cierre de emergencia con pendientes mantiene bloqueo 409.} NR-26; NP-11, NP-17 (C20.1, Alta)}{Emergencia con U1 sin retornar; acciones EMERGENCY\_CLOSE y luego intento de gestión o reencolado.}{1) Ejecutar EMERGENCY\_CLOSE. 2) Consultar pendientes visibles. 3) Intentar gestión/\allowbreak{}requeue de pedido bloqueado.}{Incidencia cierra dejando unidades sin retorno visibles y pedidos afectados bloqueados con 409 RECOVERY\_PENDING.}{Pendientes siguen listados tras EMERGENCY\_CLOSE y cualquier intento de requeue/\allowbreak{}gestión retorna 409 RECOVERY\_PENDING.}
\caso{CP-089}{\textbf{Cierre de emergencia no libera repartidor si returnedAt NULL.} NR-26; NP-11, NP-17 (C20.2, Alta)}{Ruta en emergencia cerrada con EMERGENCY\_CLOSE mientras existe RouteUnitRecovery con returnedAt IS NULL.}{1) Ejecutar EMERGENCY\_CLOSE. 2) Consultar disponibilidad de REPARTIDOR-A y reservas activas.}{El cierre EMERGENCY\_CLOSE no restablece por sí solo la disponibilidad ni libera al repartidor mientras returnedAt sea NULL.}{DriverReservation permanece activa tras EMERGENCY\_CLOSE y REPARTIDOR-A no aparece disponible para nueva asignación.}
\caso{CP-090}{\textbf{Liberación de repartidor por ADMIN respeta ausencias y 409.} NR-25, NR-26; NP-11, NP-17 (C20.3, Alta)}{Ruta INTERRUPTED con EMERGENCY\_DRIVER\_RELEASE; variante con ausencia activa; segunda llamada repetida.}{1) Ejecutar EMERGENCY\_DRIVER\_RELEASE como ADMIN. 2) Verificar DriverReservation y registro de ausencia. 3) Repetir acción.}{Libera DriverReservation respetando ausencias vigentes (sigue no disponible si ausente); al repetir rechaza 409 DRIVER\_ALREADY\_RELEASED.}{Reserva liberada una sola vez; la ausencia vigente prevalece en disponibilidad y repetición devuelve 409 DRIVER\_ALREADY\_RELEASED.}
\caso{CP-091}{\textbf{Declarar pérdida con stock \ensuremath{\geq} 1 prepara reemplazo y descuenta.} NR-27; NP-12, NP-18 (C21.1, Alta)}{Original no recuperado de ítem con Product.\allowbreak{}stock = S (\ensuremath{\geq} 1); LOSS\_DECLARE con orderId, orderItemId y requestId.}{1) Ejecutar LOSS\_DECLARE como ADMIN. 2) Consultar OrderItemLoss, Product.\allowbreak{}stock y movimientos de inventario.}{Se crea OrderItemLoss con reemplazo PREPARED, Product.\allowbreak{}stock pasa a S-1 y registra movimiento REPLACEMENT\_OUT.}{Estado PREPARED sin esperar retorno original; exactamente 1 unidad descontada (S-1) con trazabilidad REPLACEMENT\_OUT.}
\caso{CP-092}{\textbf{Reemplazo mantiene mismo pedido y no genera nuevo cobro.} NR-27; NP-12, NP-18 (C21.2, Alta)}{Pedido original con 1 PurchaseReceipt y saldo inicial de CLIENTE-A tras compra original.}{1) Consultar pedido, historial de pérdida/\allowbreak{}reemplazo, PurchaseReceipt y saldo tras LOSS\_DECLARE.}{Conserva el mismo pedido e historial de pérdida/\allowbreak{}reemplazo; no crea nuevo PurchaseReceipt ni nuevo débito.}{Conteo de PurchaseReceipt igual a 1, saldo sin cargos adicionales y mismo ID de pedido.}
\caso{CP-093}{\textbf{Declarar pérdida con stock 0 crea PENDING sin stock negativo.} NR-27; NP-18 (C22.1, Alta)}{Ítem con Product.\allowbreak{}stock === 0; solicitud LOSS\_DECLARE válida con orderId, orderItemId y requestId.}{1) Ejecutar LOSS\_DECLARE como ADMIN. 2) Consultar estado de OrderItemLoss, Product.\allowbreak{}stock y movimientos.}{Crea pérdida en estado PENDING sin descuento de stock (se mantiene en 0) ni reemplazo ficticio.}{Estado PENDING, Product.\allowbreak{}stock === 0 (sin valores negativos) y cero movimientos REPLACEMENT\_OUT.}
\caso{CP-094}{\textbf{Contacto y reprogramación tras reemplazo sin entrega hoy.} NR-27; NP-18 (C22.2, Alta)}{Pedido con reemplazo PREPARED sin entrega posible hoy; acciones CONTACT y RESCHEDULE con fecha acordada.}{1) Registrar CONTACT como ADMIN. 2) Ejecutar RESCHEDULE con fecha futura válida. 3) Consultar auditoría y pedido.}{Registra acciones CONTACT y RESCHEDULE en auditoría y actualiza fecha acordada del pedido sin crear entrega ficticia.}{Fecha acordada persistida, historial auditado de CONTACT/\allowbreak{}RESCHEDULE y sin marcas de entrega adelantada.}
\caso{CP-095}{\textbf{Preparación diferida al reponer stock y 409 NO\_FREE\_STOCK.} NR-27; NP-18 (C22.3, Alta)}{Pérdida PENDING con lossId; primero stock = 0, luego reposición a stock = S (\ensuremath{\geq} 1); acción REPLACEMENT\_PREPARE.}{1) Invocar REPLACEMENT\_PREPARE con stock 0. 2) Reponer stock a S. 3) Repetir REPLACEMENT\_PREPARE y validar.}{Con stock 0 rechaza 409 NO\_FREE\_STOCK; tras reponer a S convierte pérdida a PREPARED y descuenta 1 unidad (S-1).}{409 NO\_FREE\_STOCK sin descuento en paso 1; transición PENDING\ensuremath{\rightarrow}PREPARED y stock S-1 en paso 3.}
\caso{CP-096}{\textbf{Recuperación tardía SELLABLE suma +1 sin doble entrega.} NR-28; NP-12, NP-18 (C23.1, Alta)}{Pérdida con lossId y Product.\allowbreak{}stock = S; acción ORIGINAL\_RECOVER con recoveryCondition = SELLABLE.}{1) Ejecutar ORIGINAL\_RECOVER SELLABLE como ADMIN. 2) Consultar Product.\allowbreak{}stock, movimiento RECOVERY\_IN y pedido.}{Product.\allowbreak{}stock aumenta a S+1 con movimiento RECOVERY\_IN, conserva historial de pérdida/\allowbreak{}reemplazo y no duplica entrega.}{Stock exacto S+1, registro de condición SELLABLE, historial íntegro y pedido sin entrega duplicada.}
\caso{CP-097}{\textbf{Recuperación tardía DAMAGED registra estado sin sumar stock.} NR-28; NP-12, NP-18 (C23.2, Alta)}{Pérdida con lossId y Product.\allowbreak{}stock = S; acción ORIGINAL\_RECOVER con recoveryCondition = DAMAGED.}{1) Ejecutar ORIGINAL\_RECOVER DAMAGED como ADMIN. 2) Consultar condición en OrderItemLoss y Product.\allowbreak{}stock.}{Registra recuperación tardía con condición física DAMAGED sin aumentar Product.\allowbreak{}stock (se mantiene en S).}{Condición DAMAGED persistida en OrderItemLoss, Product.\allowbreak{}stock igual a S y sin movimiento RECOVERY\_IN vendible.}
\caso{CP-098}{\textbf{Repetición de ORIGINAL\_RECOVER rechaza 409 sin doble suma.} NR-28; NP-12, NP-18 (C23.3, Alta)}{Mismo lossId ya recuperado previamente; segundo envío de ORIGINAL\_RECOVER con nuevo requestId.}{1) Enviar ORIGINAL\_RECOVER sobre lossId ya recuperado. 2) Verificar código de respuesta, stock y movimientos.}{Operación rechazada con 409 ORIGINAL\_ALREADY\_RECOVERED sin sumar ni descontar unidades adicionales.}{HTTP 409 ORIGINAL\_ALREADY\_RECOVERED; Product.\allowbreak{}stock e historial de movimientos idénticos a antes del reintento.}
\caso{CP-099}{\textbf{Cancelación LOSS\_CANCEL restituye stock y bloquea si no aplica.} NR-27, NR-28; NP-12, NP-18 (C23.4, Alta)}{Pérdida PREPARED con stock S; pérdida recuperada/\allowbreak{}con reclamo activo o retiro en tienda; LOSS\_CANCEL.}{1) Ejecutar LOSS\_CANCEL sobre pérdida PREPARED. 2) Verificar soft-cancel y stock S+1. 3) Probar casos bloqueados.}{Soft-cancel de OrderItemLoss y restituye +1 (REPLACEMENT\_CANCEL\_IN) si estaba PREPARED; rechaza 409 en bloqueados.}{Stock S+1 solo en PREPARED, soft-cancel auditado y HTTP 409 cuando ya hubo recuperación, reclamo activo o retiro.}
\caso{CP-100}{\textbf{Concurrencia UNIT\_RECOVERY\_CONFIRM vs LOSS\_DECLARE en unidad.} NR-25 a NR-29; Transversal (C24.1, Alta)}{Misma unidad U1 pendiente en ruta INTERRUPTED; 2 hilos simultáneos: UNIT\_RECOVERY\_CONFIRM y LOSS\_DECLARE.}{1) Disparar ambas peticiones concurrentemente con requestId distintos. 2) Inspeccionar respuestas y BD.}{Exactamente una operación gana (200) y la otra rechaza con 409 UNIT\_HAS\_ACTIVE\_LOSS o 409 UNIT\_ALREADY\_RECOVERED.}{Sin estado inconsistente ni doble mutación: 1 éxito y 1 conflicto 409 específico según orden de serialización.}
\caso{CP-101}{\textbf{Idempotencia de intervenciones con requestId UUID y conflicto.} NR-25 a NR-29; Transversal (C24.2, Alta)}{Mismo requestId UUID: 1) envío válido, 2) repetición idéntica, 3) reenvío con payload o acción distinta.}{1) Ejecutar intervención válida. 2) Reenviar payload idéntico con mismo requestId. 3) Reenviar payload distinto.}{Repetición idéntica re-entrega auditoría original sin re-ejecutar efecto; payload distinto rechaza 409 ConflictException.}{1 sola mutación de estado, mismo registro de auditoría en replay y 409 ConflictException ante payload alterado.}
\caso{CP-102}{\textbf{Rendimiento generación de rutas con volumen de 100 pedidos.} NR-17 a NR-19; Plan Sec. 4 (C29.1, Complementaria)}{Dataset sintético aislado de 100 pedidos elegibles del día dentro de cobertura vial.}{1) Registrar versión y hardware del entorno. 2) Medir GET /\allowbreak{}api/\allowbreak{}v1/\allowbreak{}route-proposals con 100 pedidos.}{Se registran latencia p50, máximo, tasa de errores, versión y hardware sin clasificar contra un SLA inventado.}{Reporte reproducible con p50, max, errores y especificaciones de entorno/\allowbreak{}versión para línea base de 100 pedidos.}
\caso{CP-103}{\textbf{Rendimiento generación de rutas con volumen de 500 pedidos.} NR-17 a NR-19; Plan Sec. 4 (C29.2, Complementaria)}{Dataset sintético aislado de 500 pedidos elegibles del día con distribución horaria y vial.}{1) Mantener misma versión y hardware que en 100 pedidos. 2) Medir GET /\allowbreak{}api/\allowbreak{}v1/\allowbreak{}route-proposals con 500 pedidos.}{Se documentan p50, máximo, errores o impedimento real observado en entorno aislado sin inventar SLA fijo.}{Evidencia completa de p50, max, errores, versión y hardware para el volumen de 500 pedidos.}
\caso{CP-104}{\textbf{Rendimiento generación de rutas con volumen de 2000 pedidos.} NR-17 a NR-19; Plan Sec. 4 (C29.3, Complementaria)}{Dataset sintético aislado de 2000 pedidos elegibles del día; referencia comparativa 100/\allowbreak{}500.}{1) Medir GET /\allowbreak{}api/\allowbreak{}v1/\allowbreak{}route-proposals con 2000 pedidos. 2) Comparar tiempos y errores frente a 100 y 500 pedidos.}{Se registran p50, máximo, errores, versión y hardware junto con comparativa 100/\allowbreak{}500/\allowbreak{}2000 sin umbral inventado.}{Tabla comparativa reproducible 100/\allowbreak{}500/\allowbreak{}2000 con métricas y entorno reales, sin calificar aprobación por SLA arbitrario.}
\caso{CP-105}{\textbf{Concurrencia Shop última unidad: [201, 409 STOCK\_UNAVAILABLE].} NR-01, NR-05, NR-29; Transversal (C31.1, Alta)}{Producto con stock = 1; 2 solicitudes simultáneas válidas (CLIENTE-A y CLIENTE-B) con tarjetas demo e Idempotency-Key.}{1) Ejecutar ambos POST /\allowbreak{}api/\allowbreak{}v1/\allowbreak{}shop/\allowbreak{}checkout en paralelo. 2) Consultar respuestas, stock, pedidos y recibos.}{Respuestas [201, 409 STOCK\_UNAVAILABLE]; se crea 1 solo pedido y 1 recibo, stock final 0 sin sobreventa.}{Exactamente un 201 y un 409 STOCK\_UNAVAILABLE; stock \ensuremath{\geq} 0 (0 final) y sin débito ni pedido para la rechazada.}
\caso{CP-106}{\textbf{Concurrencia Shop saldo: [201, 422 INSUFFICIENT\_DEMO\_FUNDS].} NR-01, NR-05, NR-29; Transversal (C31.2, Alta)}{Tarjeta DEMO\_LIMITED con saldo 2000c; producto con stock \ensuremath{\geq} 2; 2 compras concurrentes de 1725c.}{1) Lanzar 2 POST /\allowbreak{}api/\allowbreak{}v1/\allowbreak{}shop/\allowbreak{}checkout simultáneos de 1725c. 2) Verificar respuestas, saldo final y stock.}{Respuestas [201, 422 INSUFFICIENT\_DEMO\_FUNDS]; saldo final 275c y stock descontado solo 1 vez.}{1 compra aprobada (201) y 1 rechazada (422 INSUFFICIENT\_DEMO\_FUNDS); saldo final 275c y stock intacto en el rechazo.}
\caso{CP-107}{\textbf{Fallo pre-commit Shop con trigger en purchaseReceipt: rollback.} NR-01, NR-05, NR-29; Transversal (C31.3, Alta)}{Prueba de integración/\allowbreak{}BD aislada: trigger BEFORE INSERT controlado en purchaseReceipt; checkout válido.}{1) Activar trigger controlado en entorno aislado. 2) Ejecutar POST /\allowbreak{}api/\allowbreak{}v1/\allowbreak{}shop/\allowbreak{}checkout. 3) Consultar BD y retirar trigger.}{Checkout falla con HTTP 500 y realiza rollback total: sin pedido, sin PurchaseReceipt, sin descuento de stock ni saldo.}{HTTP 500; conteos de Order, PurchaseReceipt, Product.\allowbreak{}stock y saldo demo idénticos al estado previo a la petición.}
\end{disenio}
